\chapter{Projective geometry}\label{geometry:ch:projective}

\begin{center}\itshape
Metrical geometry is thus a part of descriptive geometry, and descriptive geometry is all geometry.
\upshape(Arthur Cayley, \emph{A Sixth Memoir upon Quantics}, 1859)
\end{center}

A tiled floor keeps its parallel edges, but a picture of it sends them towards vanishing points. Directions lying in that one floor plane share a horizon line. Distances and right angles may change in the picture, yet lines remain lines and incidence remains recognizable. Projective coordinates let us calculate those changes without treating parallelism as an exception.

Brunelleschi's early-fifteenth-century perspective demonstration and Alberti's \emph{De pictura} of 1435 made projection a practical method for painters. Desargues developed its geometry in his 1639 work on conics. Poncelet returned to these ideas during and after his captivity at Saratov in 1812--1814, publishing his treatise in 1822. The accounts listed below distinguish that history from the coordinate proofs we will use.

The algebra comes from Volume~3: fields, bases, rank and determinants, with finite fields for the last part. A projective point will be a one-dimensional vector subspace. This makes points at infinity as ordinary as other points and lets two lines always meet. Cayley's broad claim in the epigraph concerns how metric geometries can be recovered after extra projective data are fixed; Chapter~9 gave one such connection with the hyperbolic plane.

\begin{think}
Join corresponding vertices of two triangles to a common point. Then intersect corresponding side lines. The three resulting points line up. What changes if the triangles lie in different planes in space, where their planes themselves meet in a line?
\end{think}

\section{The projective plane}

\begin{definition}\label{geometry:def:projplane}
For a field $K$, the projective plane $\mathbb P^2(K)$ is the set of one-dimensional subspaces of $K^3$. A nonzero vector $p=(x,y,z)$ represents the point $[p]=[x:y:z]$; multiplying all coordinates by a nonzero scalar leaves the point unchanged. These are \emph{homogeneous coordinates}. A projective line consists of the points represented by vectors in one fixed two-dimensional subspace of $K^3$.
\end{definition}

Such a subspace is the kernel of $ax+by+cz$ for nonzero coefficient vector $\ell=(a,b,c)$, also determined up to scale. Point $[p]$ lies on line $[\ell]$ exactly when $\ell\cdot p=0$. This coordinate pairing is bilinear and nondegenerate over any field; it is not being used as a positive Euclidean length.

Where $z\ne0$, each point has a unique representative $[x:y:1]$, giving the \emph{affine chart} $K^2$. The other points form the line $z=0$ at infinity. An affine line $ax+by+c=0$, with $(a,b)\ne(0,0)$, acquires one infinity point $[b:-a:0]$, its direction. Distinct parallel affine lines share that point.

\begin{proposition}\label{geometry:prop:incidence}
Two distinct projective points lie on exactly one line. Two distinct projective lines meet in exactly one point.
\end{proposition}
\begin{proof}
Independent representatives $p,q$ span the unique two-dimensional subspace containing their one-dimensional spans. Dually, distinct two-dimensional subspaces of $K^3$ have sum $K^3$ and intersection of dimension $2+2-3=1$. These are precisely the required line and point.
\end{proof}

\begin{theorem}[Duality]\label{geometry:thm:duality}
Sending a point $[p]$ to the line with coefficient vector $p$, and a line $[\ell]$ to the point represented by $\ell$, interchanges points and lines bijectively and preserves incidence. Consequently any incidence theorem remains true with point/line, lying on/passing through, and collinearity/concurrence exchanged, including its corresponding hypotheses.
\end{theorem}
\begin{proof}
Both incidence assertions say $\ell\cdot p=0$, since the coordinate pairing is symmetric. Apply any given incidence theorem to the interchanged configuration and translate its conclusion back. Because the interchange is its own inverse, this proves the dual statement with all its hypotheses as well.
\end{proof}

The cross product gives a convenient coordinate formula over every field. Its components are alternating $2\times2$ minors, so $(p\times q)\cdot r=\det(p,q,r)$, and it is nonzero for independent $p,q$. Thus $[p\times q]$ is their joining line; the same formula on line coefficients gives their intersection. Three represented points are collinear exactly when their $3\times3$ determinant is zero. None of this requires an ordered field or characteristic different from two.

\begin{figure}[ht]
\centering
\fig{geometry-ch04-floor}
\caption{A tiled floor seen from a height. The picture is the image of the floor under a projective map. The parallel edges of the tiles
meet at the central vanishing point, the diagonals at two others, and all three lie on the horizon, the image of the line at infinity.}
\label{geometry:fig:floor}
\end{figure}

\section{Projective transformations}

An invertible matrix on $K^3$ induces a bijection of projective points and sends lines onto lines. Only scalar matrices act trivially: fixing the three coordinate points forces the matrix to be diagonal, and fixing $[1:1:1]$ forces the three diagonal entries equal. The group of \emph{projective transformations} is therefore $\mathrm{PGL}_3(K)=\mathrm{GL}_3(K)/K^\times$.

Using $K^2$ similarly gives $\mathbb P^1(K)$, with $x=[x:1]$ and $\infty=[1:0]$. Its group acts by $(\alpha x+\beta)/(\gamma x+\delta)$, with nonzero determinant. When $\gamma\ne0$ the pole maps to infinity and infinity to $\alpha/\gamma$; when $\gamma=0$ the affine map fixes infinity. Over $\C$ these agree with Appendix~B's M\"obius maps.

For a real perspective model let eye height $h>0$ and picture distance $f>0$. A floor point $(x,y)$, with $y$ measured beyond the picture plane, appears at $(fx/(y+f),-fh/(y+f))$ relative to eye level. Homogeneously the map is
\[
[x:y:1]\longmapsto[fx:-fh:y+f],
\]
with matrix determinant $f^2h\ne0$. The floor's infinity points $[x:y:0]$ map to $[fx:0:y]$, forming the horizon at eye level. Only the floor's directions are being collected on this horizon; directions in a different spatial plane need not share it.

\begin{definition}
On $\mathbb P^1(K)$, general-position points are distinct. On $\mathbb P^2(K)$, a general-position collection has no three collinear and no repeated points. A \emph{frame} in dimension $n=1,2$ is an ordered collection of $n+2$ general-position points.
\end{definition}

\begin{proposition}[Frames]\label{geometry:prop:frames}
For $n=1$ or $2$, exactly one projective transformation carries any frame of $\mathbb P^n(K)$ to any other. In particular a matrix inducing the identity is scalar.
\end{proposition}
\begin{proof}
Start from the standard frame $[e_1],\ldots,[e_{n+1}],[e_1+\cdots+e_{n+1}]$. For a target frame, choose representatives $p_1,\ldots,p_{n+2}$. The first $n+1$ form a basis, so write $p_{n+2}=\sum_i\lambda_ip_i$. No coefficient vanishes, since that would place the last point on the span of the other $n$ representatives, violating general position. The matrix with columns $\lambda_ip_i$ is invertible and sends the standard frame to the target.

A map fixing the standard frame has a diagonal matrix from its first $n+1$ points. The final sum point makes all diagonal entries equal, so it is scalar and induces the identity. This proves uniqueness; composing two standard-frame maps gives existence between arbitrary frames.
\end{proof}

Over the reals, four general-position points can therefore be mapped to the four vertices of a square. Such a coordinate change need not preserve any Euclidean right angle. A perspective photograph can preserve the incidence of a tile grid while losing its original lengths; recovering those requires additional calibration.

\section{Desargues' theorem}

Two triangles are \emph{in perspective from a point} if the lines joining corresponding vertices are concurrent. They are \emph{in perspective from a line} if the intersections of corresponding side lines are collinear. Distinctness assumptions ensure those joins and intersections are actually defined.

\begin{theorem}[Desargues]\label{geometry:thm:desargues}
Let $ABC,A'B'C'$ be nondegenerate projective triangles with $A\ne A'$, $B\ne B'$, $C\ne C'$, and distinct corresponding side lines. If $AA',BB',CC'$ meet at a point $O$ different from all six vertices, then
$AB\cap A'B'$, $BC\cap B'C'$ and $CA\cap C'A'$ are collinear.
\end{theorem}
\begin{proof}
Write $O=[o]$. Collinearity of the three distinct points $O,A,A'$ expresses $o=\alpha a+\alpha'a'$ with both coefficients nonzero. Rescale the representatives, including a minus sign in the second, to obtain $o=a-a'$. Do this for each pair, so
\[
o=a-a'=b-b'=c-c'.
\]
Then $a-b=a'-b'$ is nonzero and lies in both side subspaces for $AB,A'B'$. Since those lines differ, it represents their unique intersection. Likewise $b-c$ and $c-a$ represent the other two intersections. Their sum is zero, making the three representatives dependent. They lie in a two-dimensional subspace, or a smaller subspace contained in one, so their projective points are collinear.
\end{proof}

Duality gives the converse with its dual nondegeneracy conditions. In particular, for the corresponding distinct vertices and sides, an axis of side intersections different from all six side lines gives concurrence at a point different from the six vertices. Degenerate coincidences require their own interpretation rather than dropping those hypotheses.

\begin{figure}[ht]
\centering
\fig{geometry-ch04-desargues}
\caption{Desargues' theorem. The triangles $ABC$ and $A'B'C'$ are in perspective from $O$; the intersections $P,Q,R$ of corresponding sides
lie on one line (rust).}
\label{geometry:fig:desargues}
\end{figure}

\begin{beauty}
In projective three-space, two distinct planes meet in a line. If triangles in those planes are perspective from an outside point, each pair of corresponding side lines meets in both planes; their three intersections therefore lie on the common line.

The normalization in the proof supplies a precise lift of our planar configuration to that spatial picture. In $K^4$, take representatives $(a,0),(b,0),(c,0)$ for the first triangle and $(a',1),(b',1),(c',1)$ for the second. Their corresponding joins pass through $(o,-1)$ because their differences equal that vector. The two triangles span distinct projective planes, and projection $(x,t)\mapsto x$ recovers the original configuration. Thus the plane-incidence statement has a spatial model whose vertices and joining lines are all specified.
\end{beauty}

\section{Pappus' theorem}

Pappus recorded this theorem about two lines in Book~VII of his \emph{Collection}, around the early fourth century. Its coordinate proof will use commutativity of the field, unlike the vector subtraction in Desargues.

\begin{theorem}[Pappus]\label{geometry:thm:pappusline}
Let distinct lines $\ell,\ell'$ meet at $O$. Choose distinct $A,B,C$ on $\ell$ and distinct $A',B',C'$ on $\ell'$, all avoiding $O$. Then
\[
X=AB'\cap A'B,\qquad Y=AC'\cap A'C,\qquad Z=BC'\cap B'C
\]
are collinear.
\end{theorem}

\begin{think}
Where does multiplying two coordinates enter? The vector relation at the end will expose exactly which products need to commute.
\end{think}

\begin{proof}
Choose representatives $u,v,o$ with $[o]=O$, $[u]\in\ell\setminus\{O\}$ and $[v]\in\ell'\setminus\{O\}$. They form a basis. In these coordinates the six points are
\[
A=[1:0:a],\ B=[1:0:b],\ C=[1:0:c],\qquad
A'=[0:1:a'],\ B'=[0:1:b'],\ C'=[0:1:c'],
\]
with the entries in each triple distinct. Joining and intersecting by cross products gives representatives
\[
\begin{aligned}
x&=(a'-b',a-b,aa'-b'b),\\
y&=(a'-c',a-c,aa'-c'c),\\
z&=(b'-c',b-c,bb'-c'c).
\end{aligned}
\]
For instance the two line vectors for $X$ are $(-a,-b',1)$ and $(-b,-a',1)$, producing the first formula. The representatives are nonzero by distinctness. Their first two coordinates cancel in $x-y+z$, while the third leaves $bb'-b'b=0$. Hence they are dependent and the three intersections are collinear.
\end{proof}

\begin{figure}[ht]
\centering
\fig{geometry-ch04-pappus}
\caption{Pappus' theorem. The three crossing points $X,Y,Z$ of the ``wrong'' pairs of lines are collinear.}
\label{geometry:fig:pappus4}
\end{figure}

\begin{beauty}
The last cancellation shows where this field proof uses commutativity. It does not license the same determinant calculation over the quaternions: determinants and left/right scalar conventions there require different arguments. We quote the broader coordinatization theorem: a Desarguesian projective plane can be coordinatized by a division ring, and Pappus holds in that plane exactly when the division ring is commutative. Hilbert and Hartshorne discuss this connection. The theorem turns an incidence property into an algebraic restriction, but its division-ring converse has not been proved by the calculation above.
\end{beauty}

\begin{remark}
Hessenberg's 1905 theorem that Pappus implies Desargues in an abstract projective plane is also quoted here. Non-Desarguesian projective planes exist. A standard example is the projective completion of Moulton's affine plane, whose negative-slope lines bend at the vertical axis. The affine construction alone is not yet a projective plane; its completion adds the appropriate infinity points. These examples and their incidence checks are in Hartshorne, and none is used as a premise for our field-coordinate proofs.
\end{remark}

\section{The cross-ratio}

A midpoint can move under projection. Four collinear points still carry a quantity that survives: first choose coordinates in which three of them become $\infty,0,1$, then read the coordinate of the fourth.

\begin{definition}\label{geometry:def:crossratio}
For $a,b,c,d\in\mathbb P^1(K)$ with $a,b,c$ distinct, let $g$ be the unique projectivity taking $a,b,c$ to $\infty,0,1$. Define their \emph{cross-ratio} by $(a,b;c,d)=g(d)\in K\cup\{\infty\}$.
\end{definition}

The fourth point may equal one of the first three: the answers are then $\infty,0,1$, respectively. If $a,b,c$ are finite, their normalizing map is $g(x)=(c-a)(x-b)/((c-b)(x-a))$, so for finite $d$,
\begin{equation}\label{geometry:eq:crossratio}
(a,b;c,d)=\frac{(c-a)(d-b)}{(c-b)(d-a)}.
\end{equation}
A zero denominator means infinity. To include infinity over any field, write $[u,v]=\det(u,v)$ for two column representatives in $K^2$. The homogeneous answer is
\[
\bigl[[c,a][d,b]:[c,b][d,a]\bigr].
\]
The first three points being distinct ensures these two entries never vanish together. Rescaling any representative multiplies both entries by the same factor. Thus this formula works in finite fields too, without an appeal to limits. In particular $(a,b;c,\infty)=(c-a)/(c-b)$ when $a,b,c$ are finite.

\begin{theorem}\label{geometry:thm:crossratio}
Every $h\in\mathrm{PGL}_2(K)$ preserves cross-ratios. Conversely, if $(a,b;c,d)=(a',b';c',d')$, with both first triples distinct, a projectivity takes $a,b,c,d$ respectively to $a',b',c',d'$.
\end{theorem}
\begin{proof}
If $g$ normalizes $a,b,c$, then $g h^{-1}$ normalizes their images under $h$. Evaluating at $h(d)$ gives $g(d)$, proving invariance. Conversely let $g,g'$ normalize the two triples. The map $h=(g')^{-1}g$ takes the triples to each other, and equality of cross-ratios gives $h(d)=d'$. This also covers repeated fourth points.
\end{proof}

On a plane line $L=[V]$, choose a linear isomorphism $K^2\to V$ and use it to define the cross-ratio. Two choices differ by an invertible $2\times2$ matrix, so invariance proves independence of the choice. A projectivity of the whole plane therefore preserves cross-ratios on every line.

\begin{proposition}\label{geometry:prop:perspectivity}
Let $L,M$ be projective lines and $O$ a point on neither. Central projection, $X\mapsto OX\cap M$, is a projectivity $L\to M$ and preserves cross-ratios.
\end{proposition}
\begin{proof}
Write $L=[V],M=[W],O=[o]$. The direct sum $K^3=W\oplus Ko$ gives a linear projection $\pi$ onto $W$. Restricted to $V$ its kernel is zero, since $o\notin V$; hence $\pi|_V:V\to W$ is an isomorphism. Also $\pi(x)=x-\lambda o$ represents $OX\cap M$, proving the claim. If $L=M$, the projection is the identity.
\end{proof}

Four equally spaced real posts with coordinates $0,1,2,3$ have cross-ratio $4/3$. Their images need not be equally spaced, but still have that cross-ratio. This gives a way to check perspective measurements without knowing the camera distance.

For four distinct points, their $24$ orderings give at most six values:
\[
\lambda,\quad\frac1\lambda,\quad1-\lambda,\quad
\frac1{1-\lambda},\quad\frac{\lambda-1}{\lambda},\quad
\frac\lambda{\lambda-1}.
\]
Some coincide; the exercise below treats the characteristic-dependent exceptions.

\begin{definition}
A collinear quadruple is \emph{harmonic} when $(A,B;C,D)=-1$. The point $D$ is then the harmonic conjugate of $C$ with respect to $A,B$.
\end{definition}

For four distinct points this requires characteristic different from two. In characteristic two, $-1=1$ forces $D=C$. For the moment assume $1+1\ne0$. With $D=\infty$, harmonicity says $(c-a)/(c-b)=-1$, or $c=(a+b)/2$. A straightedge can construct this relation through incidences alone. A \emph{complete quadrangle} consists of four general-position points and their six joining lines.

\begin{theorem}[Harmonic construction]\label{geometry:thm:harmonic}
Suppose $1+1\ne0$ in $K$. For a complete quadrangle $EFGH$, set $A=EF\cap GH$, $B=EH\cap FG$. Let $C=EG\cap AB$, $D=FH\cap AB$. Then $(A,B;C,D)=-1$; in particular $D$ is determined by $A,B,C$, independently of the auxiliary quadrangle.
\end{theorem}
\begin{proof}
Frames and cross-ratio invariance allow the coordinates
\[
E=[1:1:1],\ F=[1:-1:1],\ G=[-1:1:1],\ H=[-1:-1:1].
\]
Each triple has determinant $4$ or $-4$, hence is independent. The lines $EF,GH$ have equations $x=z,x=-z$, giving $A=[0:1:0]$; $EH,FG$ have equations $y=x,y=-x$, giving $B=[0:0:1]$. On $AB$, where $x=0$, the remaining lines $EG:y=z$ and $FH:y=-z$ give $C=[0:1:1]$, $D=[0:1:-1]$. The coordinate $[0:s:t]\mapsto[s:t]$ makes these $\infty,0,1,-1$, proving the claim. The normalizing projectivity is bijective, so exactly one $D$ can have value $-1$ for fixed distinct $A,B,C$.
\end{proof}

In characteristic two the construction collapses. For the standard frame $e_1,e_2,e_3,e_1+e_2+e_3$, the three diagonal points are $[1:1:0],[1:0:1],[0:1:1]$. Their determinant is $-2$, so they lie on a line. Frames transport this to every quadrangle. The first two are $A,B$ and the third is $EG\cap FH$, so the latter lies on $AB$ and $C=D$. This is an incidence feature of the Fano plane, not a failure of the coordinate notation.

\section{The fundamental theorem}

An invertible matrix sends lines to lines. Could a bijection preserve all lines without coming from such a matrix? Complex conjugation does. Over the reals, incidence itself will force addition, multiplication and finally every coordinate to stay fixed after a frame is normalized.

\begin{think}
Using joins, intersections and parallels, construct $x+y$ and $xy$ on a coordinate axis. A line-preserving map must respect those constructions. What freedom can its coordinate function still have?
\end{think}

\begin{theorem}[Fundamental theorem of projective geometry]\label{geometry:thm:fundamental}
A bijection $f:\mathbb P^2(\mathbb R)\to\mathbb P^2(\mathbb R)$ that maps each line onto a line is a projective transformation.
\end{theorem}
\begin{proof}
First, the inverse also sends lines onto lines. For two distinct points of a target line, join their preimages: its image is exactly that target line, and injectivity gives its full preimage. Thus collinearity is preserved in both directions and a frame goes to a frame. Compose $f$ with a projectivity to make it fix the standard frame.

The line at infinity joins two fixed coordinate points and is preserved. Its complement, the affine plane, is preserved too; parallelism is incidence on that infinity line. The origin, the point $(1,1)$ and both axis directions are fixed. Hence both axes, and the points $(1,0),(0,1)$ found by vertical and horizontal intersections, are fixed.

Define $\sigma$ by $f(x,0)=(\sigma(x),0)$. It is bijective and fixes $0,1$. The direction of the line joining $(1,0)$ and $(0,1)$ is fixed. The parallel of this direction through $(y,0)$ meets the other axis at $(0,y)$; hence $f(0,y)=(0,\sigma(y))$. Intersecting vertical and horizontal lines now gives
\[
f(x,y)=(\sigma(x),\sigma(y)).
\]

To recover addition, join $(0,1)$ to $(x,0)$ and take its parallel through $(y,1)$. Its points are $(y+tx,1-t)$, so its intersection with the horizontal axis is $(x+y,0)$. Applying $f$ to the construction gives $\sigma(x+y)=\sigma(x)+\sigma(y)$. For multiplication, the parallel through $(x,0)$ to the line from $(1,0)$ to $(0,y)$ has points $(x-t,ty)$. At the vertical axis it gives $(0,xy)$, so $\sigma(xy)=\sigma(x)\sigma(y)$. These constructions also work when either input is zero. Thus $\sigma$ is a field automorphism.

A real field automorphism preserves nonnegative numbers: $u=s^2$ implies $\sigma(u)=\sigma(s)^2\ge0$. It consequently preserves order. It fixes integers and rationals by the field operations. For rationals $p<x<q$, order gives $p\le\sigma(x)\le q$; density of the rationals forces $\sigma(x)=x$. No continuity assumption was needed.

So $f$ fixes the affine plane pointwise. Each infinity point is the unique infinity point on a line through two affine points, and is therefore fixed. The normalized $f$ is the identity, proving the theorem.
\end{proof}

\begin{remark}
Before the real-order argument, the same proof over any field gives a field automorphism $\sigma$. Infinity directions transform as $[u:v:0]\mapsto[\sigma(u):\sigma(v):0]$, by looking at the line from the origin to $(u,v)$. Thus the general line-preserving bijection is a projectivity composed with coordinatewise application of a field automorphism. Such a map is called \emph{semilinear}. Over $\mathbb Q$ and $\mathbb R$ the automorphism must be the identity; over $\mathbb C$ conjugation is one possibility; over $\mathbb F_{p^k}$ the map $x\mapsto x^p$ is another. Von Staudt's real theorem dates to his \emph{Geometrie der Lage} (1847); Hartshorne treats the general case.
\end{remark}

\section{Finite projective planes}\label{geometry:sec:finiteplanes}

A finite field gives a finite geometry with the same incidence proofs. Write $\mathbb F_q$ for a field with $q$ elements, where $q$ is a prime power at least two, as constructed in Volume~3.

\begin{proposition}\label{geometry:prop:countplane}
The plane $\mathbb P^2(\mathbb F_q)$ has $q^2+q+1$ points and the same number of lines. Each line contains $q+1$ points and each point lies on $q+1$ lines.
\end{proposition}
\begin{proof}
The $q^3-1$ nonzero vectors split into classes of $q-1$ representing the same point. Their quotient gives $q^2+q+1$. A line comes from a two-dimensional space, with $(q^2-1)/(q-1)=q+1$ points. Duality gives the counts for lines and pencils through a point.
\end{proof}

At $q=2$ this is the seven-point \emph{Fano plane}. Its points are the seven nonzero binary vectors of length three; three distinct points lie on a line exactly when their sum is zero. Fano used it in 1892 to show that an incidence axiom need not follow from the others.

\begin{figure}[ht]
\centering
\fig{geometry-ch04-fano}
\caption{The Fano plane $\mathbb P^2(\mathbb F_2)$. The seven lines are the six segments and the circle; each line has three points, and three
points form a line when their coordinates add up to $0$ modulo $2$.}
\label{geometry:fig:fano}
\end{figure}

An \emph{abstract projective plane} is a set of points and lines satisfying three axioms: two distinct points lie on exactly one line; two distinct lines meet in exactly one point; and there are four points with no three collinear. If the plane is finite, every line has the same number $n+1$ of points, every point lies on $n+1$ lines, and there are $n^2+n+1$ points and lines. The exercise below proves these counts from the axioms; $n\ge2$ is its \emph{order}.

\begin{remark}
Field planes have prime-power order, but they do not exhaust finite projective planes: Veblen and Wedderburn constructed non-Desarguesian planes of order nine in 1907. Whether every finite plane has prime-power order remains open. We quote two restrictions, neither proved here. Bruck and Ryser (1949) showed that if a plane has order $n\equiv1$ or $2\pmod4$, then $n$ is a sum of two integer squares. This excludes $6,14,21$, but not $10$. Lam, Thiel and Swiercz established nonexistence at order ten in 1989 by a computer-assisted exhaustive proof, documented in Lam's account below. Order twelve remains open. These are quoted theorems, distinct from trying examples numerically.
\end{remark}

The line sets also give useful designs: subsets of a finite symbol set with controlled intersections. Use points as symbols and lines as cards, and every pair of cards shares exactly one symbol. For coding, put the seven nonzero binary vectors into the columns of a matrix $H$. The binary words $v$ satisfying $Hv=0$ form the length-seven Hamming code. A word with three ones belongs to it exactly when the three selected points form a Fano line. A new exercise develops the error-correction rule from this equation.

\FloatBarrier
\begin{figure}[ht]
\centering
\fig{v9-appc-dual-pappus}
\caption{Two pencils at $U,V$. Write $P_1=a\cap b'$, $P_2=a'\cap b$, $Q_1=a\cap c'$, $Q_2=a'\cap c$, $R_1=b\cap c'$, $R_2=b'\cap c$. The three coral joining lines meet at $K$: the dual of Pappus.}
\label{fig:v9-dual-pappus}
\end{figure}
\begin{figure}[ht]
\centering
\fig{v9-appc-perspectivity}
\caption{Two central projections take $X\in L$ to $Y\in N$ and then to $Z\in M$. Here $L:y=0$, $M:x=0$, $N:y=x+1$, with centres $O_1=(1,1)$ and $O_2=(-1,-2)$. The composite need not fix $S=L\cap M$.}
\label{fig:v9-two-perspectivities}
\end{figure}

Lines, frames and cross-ratios now fit together: incidence determines coordinates up to a field automorphism, while a frame removes the projective freedom. Over the reals, even that field automorphism disappears. The exercises turn these proofs into straightedge constructions, camera calibration and an error-correcting code. Appendix~D keeps this language and asks what happens when lines are replaced by conics.
\FloatBarrier

\section{Exercises}

\begin{exercise}[$\star$]\label{geometry:ex:crossproduct}
Over any field, prove that $[p\times q]$ joins distinct points $[p],[q]$, and that $[\ell\times m]$ intersects distinct lines $[\ell],[m]$. Over $\mathbb R$, find the projective intersection of $2x+y=3$ and $4x+2y=1$.
\end{exercise}
\begin{exercise}[$\star$]\label{geometry:ex:parallelinf}
Prove that distinct affine lines meet at infinity precisely when they are parallel. Characterize all projectivities that map the affine plane onto itself.
\end{exercise}
\begin{exercise}[$\star$]\label{geometry:ex:sixvalues}
For four distinct points with cross-ratio $\lambda$, show that simultaneously exchanging $a,b$ and $c,d$ preserves $\lambda$. Determine the values obtained from all $24$ orderings and all possible coincidences among the six expressions in the text. Include characteristics two and three.
\end{exercise}
\begin{exercise}[$\star$]\label{geometry:ex:harmonicmidpoint}
In characteristic different from two, prove that $(a,b;c,\infty)=-1$ exactly when $c=(a+b)/2$. Over $\mathbb R$, construct a segment's midpoint using a straightedge and a given distinct parallel line, using the harmonic construction.
\end{exercise}
\begin{exercise}[$\star$]\label{geometry:ex:dualpappus}
Give the full dual statement of Pappus, including its hypotheses, and draw its two pencils and three concurrent lines.
\end{exercise}
\begin{exercise}[$\star$]\label{geometry:ex:projline}
Count the points of $\mathbb P^1(\mathbb F_q)$ and prove $|\mathrm{PGL}_2(\mathbb F_q)|=(q+1)q(q-1)$. Check the result by counting invertible matrices.
\end{exercise}
\begin{exercise}[$\star\star$]\label{geometry:ex:fano168}
Count the Fano plane's ordered frames and deduce that its projective group has order $168$. Check by counting invertible binary matrices. Prove that every line-preserving bijection of the Fano plane is projective.
\end{exercise}
\begin{exercise}[$\star$]\label{geometry:ex:fanodiag}
For $E=[1:0:0],F=[0:1:0],G=[0:0:1],H=[1:1:1]$, compute the diagonal points $EF\cap GH$, $EG\cap FH$, $EH\cap FG$. Show they are collinear exactly in characteristic two. Deduce the corresponding assertion for every quadrangle over $\mathbb F_2$ and over $\mathbb R$.
\end{exercise}
\begin{exercise}[$\star\star$]\label{geometry:ex:planeorder}
From the abstract projective-plane axioms, prove there is a point off any two given lines. If the plane is finite, prove that some integer $n\ge2$ gives $n+1$ points on every line, $n+1$ lines through every point, and $n^2+n+1$ points and lines in total.
\end{exercise}
\begin{exercise}[$\star\star$]\label{geometry:ex:dobble}
Construct a $57$-card deck using $57$ symbols, eight on each card and eight occurrences of each symbol, with every two cards sharing exactly one symbol. Why can the same projective-plane construction not give $43$ cards with seven symbols per card?
\end{exercise}
\begin{exercise}[$\star\star$]\label{geometry:ex:complexcollineation}
Show that coordinatewise complex conjugation is a bijection of $\mathbb P^2(\mathbb C)$ sending lines onto lines, yet is not a projectivity.
\end{exercise}
\begin{exercise}[$\star\star\star$]\label{geometry:ex:sylvestergallai}
Prove the Sylvester--Gallai theorem: a finite noncollinear set of points in the real affine plane determines a line containing exactly two of its points. Does the assertion hold for finite sets in $\mathbb P^2(\mathbb F_q)$ or $\mathbb P^2(\mathbb C)$?
\end{exercise}

\begin{exercise}[$\star\star$]\label{ex:v9-keystone-calibration}
A projectivity sends $(0,0),(1,0),(1,1),(0,1)$ respectively to $(0,0),(2,0),(3,2),(0,1)$. Find its homogeneous matrix, affine formula and inverse. Show its denominator stays positive on the unit square. Find the two vanishing points of the horizontal and vertical grid lines and the image of the line at infinity. Compare the image of the bottom edge's midpoint with the midpoint of its image. Finally compute the images and cross-ratio of the four bottom-edge points $0,1/3,2/3,1$.
\end{exercise}
\begin{exercise}[$\star\star$]\label{ex:v9-frobenius-plane}
Let $\mathbb F_4=\mathbb F_2[\omega]$ with $\omega^2+\omega+1=0$. Show that squaring all coordinates defines an involutory line-preserving bijection $\tau$ of $\mathbb P^2(\mathbb F_4)$ which is not projective. Count its fixed points, fixed lines and pairs of exchanged points. Count $\mathrm{PGL}_3(\mathbb F_4)$ and, using the general-field fundamental theorem proved above, the full group of line-preserving bijections.
\end{exercise}
\begin{exercise}[$\star\star$]\label{ex:v9-hamming-plane}
Let $H$ have as columns the seven distinct nonzero vectors of $\mathbb F_2^3$, and set $C=\ker H\subseteq\mathbb F_2^7$. The \emph{weight} of a word is its number of nonzero entries; the \emph{Hamming distance} is the weight of the difference. Determine all weights and their multiplicities in $C$. Give a rule using the \emph{syndrome} $Hv$ to correct any one changed bit of a codeword. Show every binary word is within distance one of exactly one codeword. Give a two-error pattern that the rule decodes incorrectly.
\end{exercise}
\begin{exercise}[$\star\star\star$]\label{ex:v9-factor-perspectivities}
A \emph{perspectivity} between lines is central projection from a point on neither. For distinct lines $L,M$ meeting at $S$, prove that a projectivity $f:L\to M$ is one perspectivity exactly when $f(S)=S$, and otherwise is a composite of two. Prove that a projectivity of a line to itself needs at most three perspectivities. A nonidentity such map needs at most two exactly when it has a fixed point. Conclude that the real projectivity $t\mapsto-1/t$ requires three. Count the identity as needing zero.
\end{exercise}

\section{Hints}

\begin{hint}{geometry:ex:crossproduct}
Use the determinant identity and a nonzero minor. For the real example, keep the third coordinate instead of dividing by it.
\end{hint}
\begin{hint}{geometry:ex:parallelinf}
The infinity point of $ax+by+c=0$ is $[b:-a:0]$. A matrix preserving $z=0$ has a restricted last row.
\end{hint}
\begin{hint}{geometry:ex:sixvalues}
The double transpositions form a four-element subgroup. The two operations on $\lambda$ are inversion and $1-\lambda$. Find fixed points of the five nonidentity operations; division by two is unavailable in characteristic two.
\end{hint}
\begin{hint}{geometry:ex:harmonicmidpoint}
Use the supplied parallel as the side $FH$ of the quadrangle. Its intersection with $AB$ is needed in the proof, not as a point to draw.
\end{hint}
\begin{hint}{geometry:ex:dualpappus}
Replace the two base lines by two pencil centres, and each join in Pappus by an intersection. Figure~\ref{fig:v9-dual-pappus} supplies a drawing.
\end{hint}
\begin{hint}{geometry:ex:projline}
Count scalar classes of vectors. A projectivity is uniquely determined by images of three distinct points.
\end{hint}
\begin{hint}{geometry:ex:fano168}
Choose three noncollinear points, then avoid their three sides. A map fixing the standard frame fixes the remaining third points on its sides.
\end{hint}
\begin{hint}{geometry:ex:fanodiag}
The three representatives have two entries equal to one. Their determinant is $-2$.
\end{hint}
\begin{hint}{geometry:ex:planeorder}
Use a quadrangle to find a point off both lines. Projection gives a bijection of the lines; a pencil can be counted against a line off its centre.
\end{hint}
\begin{hint}{geometry:ex:dobble}
Take field-plane lines as cards. For the obstruction, identify the required order and use the quoted Bruck--Ryser theorem.
\end{hint}
\begin{hint}{geometry:ex:complexcollineation}
Conjugate the line coefficients too. Then compare the standard frame with $[1:i:0]$.
\end{hint}
\begin{hint}{geometry:ex:sylvestergallai}
Minimize the positive distance from a point to a line determined by other points. For a complex counterexample try three triples on the coordinate lines with coordinates involving a cube root of unity.
\end{hint}
\begin{hint}{ex:v9-keystone-calibration}
Use columns for the images of the affine origin and the two infinity directions. Alternatively substitute the four corner conditions into a general matrix and solve up to scale.
\end{hint}
\begin{hint}{ex:v9-frobenius-plane}
Normalize a nonzero coordinate to one before asking whether a point is fixed. A matrix is counted by choosing independent columns in succession.
\end{hint}
\begin{hint}{ex:v9-hamming-plane}
Weight-three words correspond to Fano lines. The all-ones word lies in the kernel. A one-bit error adds exactly one column to the syndrome.
\end{hint}
\begin{hint}{ex:v9-factor-perspectivities}
A single projection fixes the intersection of its source and target lines. For the converse, join two corresponding pairs. To factor a general map, choose an auxiliary line through a suitable point $T\in M$ and first project $f^{-1}(T)$ to $T$.
\end{hint}

\section{Solutions}

\begin{solution}{geometry:ex:crossproduct}
\idea Nonzero minors give the coefficient vector of the desired line or point.
\writeup The identity $(p\times q)\cdot r=\det(p,q,r)$ vanishes at $r=p,q$. Independence of $p,q$ makes some minor, hence $p\times q$, nonzero. Its kernel is the unique line through the points. The same calculation with line coefficients proves the intersection formula, since the resulting point satisfies both equations. Over $\mathbb R$,
$(2,1,-3)\times(4,2,-1)=(5,-10,0)$, giving $[1:-2:0]$. This numerical example uses the real field: in characteristic five the two coefficient triples would instead be proportional.
\end{solution}
\begin{solution}{geometry:ex:parallelinf}
\idea Infinity points record affine directions; preserving the affine chart means preserving its complement.
\writeup The point $[b:-a:0]$ of a line is its direction. Two distinct affine lines share it exactly when their direction vectors are proportional, that is, when they are parallel. A projectivity preserves the affine plane onto itself exactly when it preserves $z=0$. Its last matrix row is then $(0,0,c)$ with $c\ne0$. Divide the matrix by $c$ to get an affine formula $v\mapsto Av+b$ with $A$ invertible. Conversely every such matrix preserves the chart, giving exactly the affine maps of Appendix~A.
\end{solution}
\begin{solution}{geometry:ex:sixvalues}
\idea Reduce permutations to six fractional linear operations, then compute their stabilizers.
\writeup Swapping the first pair or the last pair individually replaces $\lambda$ by $1/\lambda$, so swapping both preserves it. Substitution also shows that $(ac)(bd)$ preserves it. The identity and the three double transpositions form a subgroup of order four, leaving at most six results from the $24$ permutations. Swapping $b,c$ gives $1-\lambda$, by expanding the numerator of $1-\lambda$. Inversion and this operation generate precisely the six displayed expressions, so all those expressions occur.

Distinctness gives $\lambda\notin\{0,1,\infty\}$. Two expressions coincide exactly when a nonidentity one fixes $\lambda$, after composing with the inverse of one of them. The three involutions give $\lambda=-1,1/2,2$ when the characteristic is not two; none has an allowed fixed point in characteristic two. The two order-three operations have the fixed-point equation $\lambda^2-\lambda+1=0$.

If the characteristic is neither two nor three, the harmonic set $\{-1,1/2,2\}$ has an orbit of size three. Each root of $\lambda^2-\lambda+1$, when present in $K$, has an orbit of size two; these roots are not harmonic. Every other orbit has size six. In characteristic two, the roots of $\lambda^2+\lambda+1$ give size two, and all other allowed values give size six. In characteristic three, $-1=1/2=2$ and the quadratic is $(\lambda+1)^2$; the value $2$ is fixed by all six operations and has orbit size one. All other allowed values have size six. These cases exhaust the five nonidentity operations and therefore all coincidences.
\end{solution}
\begin{solution}{geometry:ex:harmonicmidpoint}
\idea Arrange the given parallel as the side whose intersection is at infinity.
\writeup The formula $(a,b;c,\infty)=(c-a)/(c-b)$ equals $-1$ precisely when $2c=a+b$. For the real construction choose $E$ off both $AB$ and the supplied parallel $m$. Set $F=EA\cap m$, $H=EB\cap m$, $G=AH\cap BF$, then $C=EG\cap AB$. Choose $E$ generically so these auxiliary intersections are finite and form a quadrangle; for example avoid the single height making $AH$ and $BF$ parallel. Here $EF\cap GH=A$, $EH\cap FG=B$, and $FH=m$. The harmonic theorem gives $(A,B;C,D)=-1$ with $D=m\cap AB=\infty$, so $C$ is the midpoint. Every drawn join and intersection is finite; $D$ appears only in the proof.
\end{solution}
\begin{solution}{geometry:ex:dualpappus}
\idea Incidence duality turns the Pappus line into a concurrence point.
\writeup Let $U,V$ be distinct points, with three distinct lines $a,b,c$ through $U$ and three distinct lines $a',b',c'$ through $V$, none equal to $UV$. Set
\[
\begin{array}{lll}
P_1=a\cap b',&Q_1=a\cap c',&R_1=b\cap c',\\
P_2=a'\cap b,&Q_2=a'\cap c,&R_2=b'\cap c.
\end{array}
\]
Then $P_1P_2,Q_1Q_2,R_1R_2$ are concurrent, by applying the duality theorem to Pappus with its distinctness hypotheses. Figure~\ref{fig:v9-dual-pappus} draws this configuration. One exact real realization has $U=(-2,0),V=(2,0)$, slopes $1/2,3/2,1$ for the first pencil and $-6/5,-1/2,-5/2$ for the second. The intersections are
\[
\begin{aligned}
P_1&=(0,1),&P_2&=(-2/9,8/3),\\
Q_1&=(4/3,5/3),&Q_2&=(2/11,24/11),\\
R_1&=(1/2,15/4),&R_2&=(-2/3,4/3).
\end{aligned}
\]
Their three joining lines pass through $K=(-12/67,157/67)$. These coordinates also provide a drawing recipe without assuming any theorem about conics.
\end{solution}
\begin{solution}{geometry:ex:projline}
\idea A projectivity acts freely and transitively on ordered triples of distinct points.
\writeup Scalar classes give $(q^2-1)/(q-1)=q+1$ points. The frames proposition identifies a projectivity uniquely with the images of a fixed triple, giving $(q+1)q(q-1)$ possibilities. Independently, an invertible matrix has $q^2-1$ choices for its first column and $q^2-q$ for its second. Dividing by the $q-1$ nonzero scalar multiples yields $(q+1)q(q-1)$ again.
\end{solution}
\begin{solution}{geometry:ex:fano168}
\idea The seven-point plane has very little freedom after a frame is chosen.
\writeup Choose the first two points in $7\cdot6$ ways. Their line has three points, leaving four choices for the third. The three sides of this triangle contain six points: its three vertices and one new point on each side. Exactly one choice remains for the fourth frame point. Hence the frame count is $7\cdot6\cdot4=168$, equal to the projective group order by simple transitivity. The matrix count agrees: $(8-1)(8-2)(8-4)=168$, with no nontrivial scalar to divide out.

For an arbitrary line-preserving bijection, frame images are frames, since inverse images of lines are lines as in the fundamental theorem. Compose with a projectivity to fix the standard frame. Each of the other three points is the third point of a line through two fixed coordinate points, hence is fixed too. The normalized map is the identity, so the original is projective.
\end{solution}
\begin{solution}{geometry:ex:fanodiag}
\idea A single determinant tests the diagonal line.
\writeup Cross products give $[1:1:0],[1:0:1],[0:1:1]$. The determinant of these three rows is $-2$, which vanishes precisely in characteristic two. Every quadrangle is the image of the standard frame under a projectivity, so the result holds for every quadrangle: its diagonal points are collinear in the Fano plane and never collinear over $\mathbb R$.
\end{solution}
\begin{solution}{geometry:ex:planeorder}
\idea Projection compares lines; a pencil then partitions all remaining points.
\writeup Choose a quadrangle $E,F,G,H$. For $L=M$, some quadrangle vertex is off that line. For distinct $L,M$, either a vertex is already off their union, or exactly two vertices lie on each, with none at $L\cap M$. Rename them $E,F\in L$, $G,H\in M$. Then $X=EG\cap FH$ lies on neither: if it lay on $L$, it would equal $E$, forcing $E,F,H$ collinear; the argument for $M$ is the same. Thus a projection centre off both exists. The maps $P\mapsto OP\cap M$ and $Q\mapsto OQ\cap L$ are inverse bijections.

In the finite case let the common line size be $n+1$. The line $EF$ already has a third point $EF\cap GH$, distinct from $E,F$, so $n\ge2$. Every point lies off some line: the three side lines of a noncollinear triangle have no common point. Projection from a point to a line off it identifies its pencil with that line, giving $n+1$ lines through it. Fixing a point $P$ partitions the other points into these $n+1$ lines minus $P$, each containing $n$ points. The total is $1+n(n+1)$. Finally, counting incidences in both ways gives the same number of lines and points, since each contributes $n+1$ incidences.
\end{solution}
\begin{solution}{geometry:ex:dobble}
\idea Point-line incidence supplies the matching-symbol rule.
\writeup Take the $57$ points of $\mathbb P^2(\mathbb F_7)$ as symbols and its $57$ lines as cards. Each line has eight points, each point lies on eight lines, and distinct lines meet once. This gives every requested count. Seven symbols per card in the same construction would require order six and $43$ points. The quoted Bruck--Ryser theorem excludes this because $6\equiv2\pmod4$ and six is not a sum of two integer squares. This excludes that full projective-plane deck; it does not prohibit smaller decks obtained by deleting cards.
\end{solution}
\begin{solution}{geometry:ex:complexcollineation}
\idea Conjugate coefficients preserve line equations, but a fixed frame detects nonprojectivity.
\writeup Conjugating $az_0+bz_1+cz_2=0$ gives the line with coefficients $\bar a,\bar b,\bar c$. Conjugation is its own inverse, so the map is bijective and onto each image line. It fixes all four standard frame points. A projectivity doing that would be the identity, but conjugation moves $[1:i:0]$ to the distinct point $[1:-i:0]$. Hence it is not projective.
\end{solution}
\begin{solution}{geometry:ex:sylvestergallai}
\idea Minimize a positive point-line distance; three points on the chosen line would yield a smaller one.
\writeup There are finitely many pairs $(P,\ell)$ with $P\in S$, with $\ell$ determined by two points of $S$, and with $P\notin\ell$. Noncollinearity makes this collection nonempty. Choose a pair of least positive distance $\delta$, and let $F$ be the perpendicular foot. If $\ell$ contained at least three points, two of them, $Q,R$, would lie on the same closed ray from $F$, with $Q$ nearer. The line $PR$ misses $Q$. Comparing triangle areas gives
\[
\operatorname{dist}(Q,PR)=\delta\frac{|QR|}{|PR|}<\delta,
\]
because $|QR|\le|FR|<|PR|$. This contradicts minimality. Thus $\ell$ contains exactly two points.

Over $\mathbb F_q$, take every projective point: every line has $q+1\ge3$ points. For a complex counterexample let $\omega$ be a primitive cube root of unity and take the nine points
\[
P_i=[1:-\omega^i:0],\quad Q_j=[1:0:-\omega^j],\quad
R_k=[0:1:-\omega^k]\quad(i,j,k=0,1,2).
\]
They are distinct and not all collinear. Pairs in one family lie on its coordinate line, containing its third point. The line through $P_i,Q_j$ has coefficients $(\omega^{i+j},\omega^j,\omega^i)$ and also contains $R_{j-i}$, with indices modulo three. Conversely any $P_i,R_k$ lies with $Q_{i+k}$ on such a line, and any $Q_j,R_k$ lies with $P_{j-k}$. Thus every determined line contains at least three of the nine points. No result about cubic curves is needed for this counterexample.
\end{solution}

\begin{solution}{ex:v9-keystone-calibration}
\idea Four corner images determine a frame map; homogeneous columns make the remaining calculations linear.
\writeup Write a matrix with columns $u,v,w$. The origin image gives $w=(0,0,c)$ with $c\ne0$. The other corner images impose $u+w\sim(2,0,1)$, $v+w\sim(0,1,1)$ and $u+v+w\sim(3,2,1)$. Solving gives, up to scale,
\[
H=\begin{pmatrix}6&0&0\\0&4&0\\-2&-1&5\end{pmatrix},\qquad
J=\begin{pmatrix}10&0&0\\0&15&0\\4&3&12\end{pmatrix}.
\]
Here $\det H=120\ne0$ and $HJ=JH=60I$. Thus
\[
(X,Y)=\frac{(6x,4y)}{5-2x-y},\qquad
(x,y)=\frac{(10X,15Y)}{4X+3Y+12}.
\]
The first denominator is at least two on the unit square. The horizontal and vertical infinity points map to $[-3:0:1]$ and $[0:-4:1]$; their joining line is $4X+3Y+12Z=0$, the image of infinity. The bottom midpoint maps to $(3/4,0)$, whereas the image edge midpoint is $(1,0)$. The four bottom coordinates map to $0,6/13,12/11,2$. Direct substitution gives
\[
\frac{(12/11)(2-6/13)}{(12/11-6/13)2}=\frac43,
\]
the same value as $(0,1/3;2/3,1)$. Cross-ratio survives although midpoint and equal spacing do not.
\end{solution}
\begin{solution}{ex:v9-frobenius-plane}
\idea Squaring preserves field operations but is not the identity; normalized coordinates separate field action from scalar freedom.
\writeup In $\mathbb F_4$, $x^4=x$, so squaring is an automorphism of order two. Squaring coordinates and line coefficients preserves incidence, and applying it twice is the identity. The standard frame is fixed, but $[\omega:1:0]$ goes to $[\omega^2:1:0]\ne[\omega:1:0]$. Thus $\tau$ cannot be projective.

Normalize the first nonzero coordinate of a fixed point to one. After squaring it remains one, so every other coordinate must satisfy $x^2=x$, giving $x\in\mathbb F_2$. There are seven fixed points, forming an embedded Fano plane. The same argument for coefficient vectors gives seven fixed lines. There are $4^2+4+1=21$ points in all, so the other fourteen form seven exchanged pairs.

Independent column choices and division by scalars give
\[
|\mathrm{PGL}_3(\mathbb F_q)|
=\frac{(q^3-1)(q^3-q)(q^3-q^2)}{q-1}.
\]
At $q=4$ this is $63\cdot60\cdot48/3=60\,480$. An automorphism of $\mathbb F_4$ fixes $0,1$ and must take $\omega$ to one of the two roots $\omega,\omega^2$ of its polynomial; both choices occur. The general-field fundamental theorem therefore gives every collineation as $g$ or $g\tau$ with $g$ projective. These two sets are disjoint, since $\tau$ is not projective, and each has $60\,480$ elements. The full order is $120\,960$.
\end{solution}
\begin{solution}{ex:v9-hamming-plane}
\idea The Fano line relation is exactly the three-bit parity relation.
\writeup The columns include a basis, so $H$ has rank three and $C$ has dimension four: there are sixteen codewords. None has weight one or two, since columns are nonzero and distinct. A weight-three word selects columns summing to zero, precisely a Fano line, giving seven words. The all-ones word belongs to $C$, since each row of $H$ contains four ones. Adding it complements every word; the seven weight-three words thus give seven weight-four words, alongside zero and the weight-seven word. These sixteen exhaust $C$. The nonzero minimum weight, hence minimum distance, is three.

If $v=c+e_i$ has a single error, then $Hv=He_i$ is column $i$. Find that column and flip bit $i$ to recover $c$. If $Hv=0$, $v$ itself is a codeword. Distinct codewords have disjoint distance-one balls, since a common word would put them within distance two. Each ball has $1+7=8$ words, and $16\cdot8=128=|\mathbb F_2^7|$, so the balls partition all words. The guarantee cannot extend to two errors: start at zero and change bits $i,j$. Their column sum is a third nonzero column $k$, so the rule flips bit $k$ and returns the nonzero codeword $e_i+e_j+e_k$ rather than zero.
\end{solution}
\begin{solution}{ex:v9-factor-perspectivities}
\idea A single projection fixes the common point; an auxiliary line can arrange that property for the second factor.
\writeup Every projection $L\to M$ fixes $S=L\cap M$. Conversely suppose $f(S)=S$, and choose distinct $P,Q\in L\setminus\{S\}$. Their images $P',Q'$ are distinct and off $S$. The lines $PP'$ and $QQ'$ meet at a point $O$ off both $L,M$: an intersection on $L$ would have to equal both $P,Q$, and similarly on $M$. Projection from $O$ agrees with $f$ at $S,P,Q$, hence agrees everywhere by the frames proposition. This proves the equivalence.

Now suppose $f(S)\ne S$. Choose $T\in M$ different from $S,f(S)$; every field line has at least three points. Put $A=f^{-1}(T)\ne S$. Choose a line $N$ through $T$ different from $M$ and $AT$, which is possible since a pencil has at least three lines. Choose $O_1\in AT\setminus\{A,T\}$. It lies off $L,N$, so its projection $\pi_1:L\to N$ takes $A$ to $T$. The projectivity $h=f\pi_1^{-1}:N\to M$ fixes $T=N\cap M$ and is therefore one projection by the first part. Thus $f=h\pi_1$ is two, and cannot be one. Figure~\ref{fig:v9-two-perspectivities} shows the construction's mechanism in real coordinates.

For a self-projectivity $f$ of $L$, first choose any projection $\pi:L\to M$ to a distinct line. The map $f\pi^{-1}:M\to L$ is at most two projections by what we just proved, giving at most three for $f$. If $f$ has a fixed point $U$, choose $M$ through $U$; then both $\pi$ and $f\pi^{-1}$ fix $U$, so two suffice. Conversely a composite $L\to N\to L$ with $N\ne L$ fixes $L\cap N$. A projection $L\to L$ is the identity, so allowing $N=L$ adds no fixed-point-free exception. A nonidentity self-projectivity cannot be one projection. Finally $t=-1/t$ has no real solution, and infinity maps to zero, so $t\mapsto-1/t$ has no projective fixed point and needs exactly three.
\end{solution}

\section*{Sources and further reading}

For incidence proofs and their axioms, Hartshorne is a useful next book; Coxeter gives many constructions to draw. Richter-Gebert develops the coordinate calculations into a broad geometric toolkit. Field and Kemp place perspective in its artistic setting. The original papers below identify the historical claims and the quoted finite-plane obstructions; those deeper results are not needed to verify our elementary coordinate proofs.

\begin{itemize}
\item A.~Cayley, \emph{A sixth memoir upon quantics}, Philosophical Transactions of the Royal Society of London 149 (1859), 61--90, \S230 (the
epigraph, shortened at the end; the full sentence continues ``\dots and reciprocally'').
\item A.~Jones, \emph{Pappus of Alexandria: Book 7 of the Collection}, Springer, 1986 (Pappus' theorem among the lemmas to Euclid's
\emph{Porisms}).
\item J.~V.~Field, \emph{The Invention of Infinity: Mathematics and Art in the Renaissance}, Oxford University Press, 1997 (Brunelleschi's
demonstration as described by Manetti, Alberti's \emph{De pictura} of 1435, and Desargues); M.~Kemp, \emph{The Science of Art}, Yale University
Press, 1990.
\item J.~V.~Field and J.~J.~Gray, \emph{The Geometrical Work of Girard Desargues}, Springer, 1987 (translation of the \emph{Brouillon project} of
1639, and Desargues' theorem as published in A.~Bosse, \emph{Mani\`ere universelle de M.~Desargues pour pratiquer la perspective}, 1648); the
\emph{Dictionary of Scientific Biography} entry on Desargues by R.~Taton (the fifty copies of 1639).
\item J.~J.~Gray, \emph{Worlds Out of Nothing: A Course in the History of Geometry in the 19th Century}, 2nd ed., Springer, 2010 (Poncelet's captivity at
Saratov and the \emph{Trait\'e} of 1822, von Staudt, Cayley's absolute).
\item R.~Hartshorne, \emph{Foundations of Projective Geometry}, Benjamin, 1967 (axiomatic planes, Desargues and Pappus, coordinatising by a skew
field, Hessenberg's theorem, the Moulton plane, the fundamental theorem); D.~Hilbert, \emph{Grundlagen der Geometrie}, Teubner, 1899;
G.~Hessenberg, \emph{Beweis des Desarguesschen Satzes aus dem Pascalschen}, Mathematische Annalen 61 (1905), 161--172; F.~R.~Moulton, \emph{A simple
non-Desarguesian plane geometry}, Transactions of the AMS 3 (1902), 192--195.
\item G.~Fano, \emph{Sui postulati fondamentali della geometria proiettiva}, Giornale di Matematiche 30 (1892), 106--132; O.~Veblen and
J.~H.~M.~Wedderburn, \emph{Non-Desarguesian and non-Pascalian geometries}, Transactions of the AMS 8 (1907), 379--388.
\item R.~H.~Bruck and H.~J.~Ryser, \emph{The nonexistence of certain finite projective planes}, Canadian Journal of Mathematics 1 (1949), 88--93;
C.~W.~H.~Lam, \emph{The search for a finite projective plane of order 10}, American Mathematical Monthly 98 (1991), 305--318.
\item H.~S.~M.~Coxeter, \emph{A problem of collinear points}, American Mathematical Monthly 55 (1948), 26--28 (Kelly's proof of the
Sylvester--Gallai theorem); M.~Aigner and G.~M.~Ziegler, \emph{Proofs from THE BOOK}, Springer, chapter ``Lines in the plane and decompositions of
graphs''.
\item H.~S.~M.~Coxeter, \emph{Projective Geometry}, 2nd ed., Springer, 1987; J.~Richter-Gebert, \emph{Perspectives on Projective Geometry}, Springer,
2011 (a modern book full of pictures and computations with cross products).
\end{itemize}

